\documentclass{article}
\usepackage[UTF8]{ctex}
\usepackage{amsmath,amssymb,amsfonts}
\usepackage{enumitem}
\usepackage{geometry}
\geometry{a4paper, margin=1in}
% === 错题本排版优化（不参与题目标识） ===
\setlength{\parindent}{0pt}
\setlength{\parskip}{0.35em plus 0.10em minus 0.08em}
\setlength{\emergencystretch}{3em}
\setlist{itemsep=0.25em, topsep=0.40em, parsep=0pt}
\allowdisplaybreaks[2]
% === 错题本排版优化结束 ===
\begin{document}
\section*{西安交通大学2024年数学分析真题}
\section*{179．西安交通大学2024年数学分析试卷（回忆版）}

\subsection{题目1}
1．求 \(\lim _{x \rightarrow 0} \frac{e^{x} \sin x-x(1+x)}{x^{3}}=\)

\subsection{题目2}
2 ．求 \(\lim _{n \rightarrow \infty}(\ln n-[\ln n])=\) \(\_\_\_\_\)．

\subsection{题目3}
3．已知 \(f(x)=x+\ln x, g(x)\) 是 \(f(x)\) 的反函数，求 \(g^{\prime \prime}(x)=\) \(\_\_\_\_\)．

\subsection{题目4}
4．求 \(\sum_{n=1}^{\infty} \frac{\sin n x}{n!}=\) \(\_\_\_\_\)．

\subsection{题目5}
5．已知 \(f(x)=|x|^{\alpha}\) ，求 \(\frac{\partial^{2} f}{\partial x_{1}^{2}}+\frac{\partial^{2} f}{\partial x_{2}^{2}}+\cdots+\frac{\partial^{2} f}{\partial x_{n}^{2}}=\) \(\_\_\_\_\)．

\subsection{题目6}
6．已知 \(f(x, y)=\left\{\begin{array}{ll}\frac{x^{3}}{x^{2}+y^{2}}, & x^{2}+y^{2} \neq 0 ; \\ 0, & x^{2}+y^{2}=0 .\end{array}\right.\) 求在 \((0,0)\) 点沿 \((\cos \alpha, \sin \alpha)\) 方向的方向导数 \(=\) \(\_\_\_\_\)．

\subsection{题目7}
7．求不定积分 \(\int \sin (\ln x) \mathrm{d} x=\) \(\_\_\_\_\)．

\subsection{题目8}
8．\(x^{2}+y^{2}=2 a z\) 和 \(x^{2}+x y+y^{2}=a^{2}\) 的交线的最大值为 \(\_\_\_\_\)．

\subsection{题目9}
9．已知 \(\Omega\) 是由

\[
(0,0,1),(0,1,1),(1,1,1),(0,0,2),(0,2,2),(2,2,2)
\]

组成的 \(\mathbb{R}^{3}\) 的一个棱台，那么 \(\iiint_{\Omega} \frac{1}{y^{2}+z^{2}} \mathrm{~d} x \mathrm{~d} y \mathrm{~d} z=\)

\subsection{题目10}
10．已知 \(S: \frac{x^{2}}{a^{2}}+\frac{y^{2}}{b^{2}}+\frac{z^{2}}{c^{2}}=1\) ，则 \(\iint_{S} x \mathrm{~d} y \mathrm{~d} z+y \mathrm{~d} z \mathrm{~d} x+z \mathrm{~d} x \mathrm{~d} y=\) \(\_\_\_\_\)．

\subsection{题目11}
11．已知数列 \(\left\{x_{n}\right\}\) 满足 \(x_{m n} \leq x_{m}+x_{n}, x_{n}>0\) ，证明： \(\lim _{n \rightarrow \infty} \sqrt[n]{x_{n}}\) 存在．

\subsection{题目12}
12．证明： \(\lim _{h \rightarrow 0^{+}, k \rightarrow 0^{+}} \frac{f\left(x_{0}+h\right)-f\left(x_{0}-k\right)}{k+h}\) 存在的充要条件为 \(f\) 在 \(x_{0}\) 处可导．

\subsection{题目13}
13．证明梅尔滕定理：设级数 \(\sum_{n=1}^{\infty} a_{n}\) 绝对收敛，\(\sum_{n=1}^{\infty} b_{n}\) 收敛，且 \(\sum_{n=1}^{\infty} a_{n}=A, \sum_{n=1}^{\infty} b_{n}=B\) ，则 \(\sum_{n=1}^{\infty} c_{n}=A B\) ，其中 \(c_{n}=a_{1} b_{n}+a_{2} b_{n-1}+\cdots+a_{n} b_{1}=\sum_{k=1}^{n} a_{k} b_{n-k+1}\) ．

\subsection{题目14}
14．已知 \(f: \mathbb{R}^{n} \rightarrow \mathbb{R}^{m}\) 是一个 \(C^{1}\) 类映射，且满足 \(f^{\prime}(\mathbf{x})=0\) ，证明：\(f\) 为常值映射．

\subsection{题目15}
15．证明：若 \(f(x)\) 在 \([a, b]\) 上连续，那么 \(f\) 在 \([a, b]\) 上有界．

\subsection{题目16}
16．设 \(a, b, T\) 是三个正实数，\(a \neq b\) ．证明：

\[
\frac{1}{T} \int_{-T}^{T} \frac{\sin a t \cos b t}{t} \mathrm{~d} t= \begin{cases}1+O\left(\frac{1}{T(a-b)}\right), & a>b \\ O\left(\frac{1}{T(a-b)}\right), & a<b\end{cases}
\]
\end{document}
